\documentclass[12pt]{article}

\usepackage[utf8]{inputenc}
\usepackage[spanish, es-tabla]{babel}
\usepackage{amsmath, amssymb}
\usepackage{graphicx}
\usepackage{booktabs}
\usepackage{geometry}
\geometry{margin=2cm}
\usepackage{hyperref}
\usepackage{listings}
\usepackage{xcolor}
\usepackage{float}
\usepackage{caption}
\usepackage{subcaption}

\hypersetup{
    colorlinks=true,
    linkcolor=blue,
    citecolor=blue,
    filecolor=magenta,
    urlcolor=cyan,
    pdftitle={Evaluación de Severidad de Quema en el Parque Nacional Blue Mountains (Black Summer 2019-2020) mediante Sentinel-2 y QGIS}
}

% Configuración para código y fórmulas lógicas
\lstset{
    breaklines=true,
    basicstyle=\ttfamily\footnotesize,
    frame=single,
    backgroundcolor=\color{gray!10},
    keywordstyle=\color{blue},
    commentstyle=\color{gray},
    stringstyle=\color{teal}
}

\title{Clasificación Supervisada del Evento Black Summer en Australia mediante Arquitecturas de Machine Learning}
\author{
  \small
  Sara Castro, Maria Alejandra Garcia, Juan David Ortiz, Eliana Pardo, Ana Maria Romero
}
\date{\today}

\begin{document}

\maketitle

\section{Introducción y Contextualización del Evento}
Durante el periodo estival 2019--2020, el este y sureste de Australia sufrieron el colapso ecológico y forestal más crítico documentado en la historia moderna del continente, conocido internacionalmente como el evento de megaincendios del ``Black Summer'' \cite{boM2020blacksummer}. Esta catástrofe fue propiciada por un clima extremo, caracterizado por una sequía plurianual sostenida y condiciones fuertemente positivas del Dipolo del Océano Índico (IOD) sumadas a fases anómalas del Modo Anular del Sur (SAM), lo cual redujo drásticamente el contenido de humedad del suelo y de la biomasa leñosa en la Gran Cordillera Divisoria \cite{abatzoglou2021anthropogenic, vanoldenborgh2021pathways}. En total, más de 24 millones de hectáreas fueron consumidas en todo el país, emitiendo pirocúmulos de humo que alcanzaron la estratosfera e impactando críticamente la biodiversidad y cuencas hidrográficas \cite{filkov2020impact, bowman2021pyrogeographic}.

Dentro del estado de Nueva Gales del Sur (NSW), uno de los eventos individuales más desastrosos fue el megaincendio de Gospers Mountain, iniciado por la caída de un rayo el 26 de octubre de 2019. Este fuego avanzó descontroladamente a través de gargantas y macizos escarpados, llegando a consumir más de 510.000 hectáreas y arrasando aproximadamente el 80\% de la superficie protegida del Área de Patrimonio Mundial de las Blue Mountains \cite{hamilton2020ecological}. La magnitud del fuego fue ampliamente documentada por medios de comunicación y comisiones de investigación (Figura \ref{fig:impacto_terreno_gospers}), evidenciando en tierra el comportamiento extremo de las llamas y la calcinación directa del estrato arbóreo frente a los equipos de emergencia.

\begin{figure}[H]
    \centering
    \includegraphics[width=0.88\textwidth]{38.-black-summer-img.jpg}
    \caption{Black Summer en Nueva Gales del Sur. Fuente: ABC News Australia \cite{abc2023gospers}.}
    \label{fig:impacto_terreno_gospers}
\end{figure}

En este contexto, la teledetección espacial y los Sistemas de Información Geográfica (SIG) constituyen herramientas indispensables para cuantificar y caracterizar espacialmente las perturbaciones por fuego sobre los ecosistemas forestales \cite{lillesand2015remote}. Con esto en mente, el propósito fundamental de este proyecto consiste en clasificar de forma supervisada las coberturas terrestres del Parque Nacional Blue Mountains (NSW) en la escena posterior al megaincendio de Gospers Mountain, y determinar cuál de cinco arquitecturas de \textit{Machine Learning} (Regresión Logística, Árboles de Decisión, Máquinas de Soporte Vectorial, Redes Neuronales Artificiales y K-Vecinos Más Cercanos) distingue mejor las coberturas a partir de su firma espectral. Para ello se emplea una imagen Sentinel-2 calibrada a nivel de reflectancia de fondo de atmósfera (BOA), de la cual se extraen las firmas espectrales de píxeles pertenecientes a \textbf{cuatro macroclases}: (1)~Agua, (2)~Urbanización, (3)~Vegetación boscosa sana y (4)~Vegetación boscosa quemada. Esta definición final de macroclases fue adoptada por recomendación del docente en la primera entrega del proyecto: se decidió \textbf{eliminar la clase Sombras} por su escasa representatividad espacial, \textbf{combinar Pasto con Vegetación boscosa sana} al compartir una firma espectral muy similar en la imagen post-incendio, y \textbf{combinar Nubes con Urbanización} dado que ambas coberturas presentan alta reflectancia en el espectro visible y resultaban indistinguibles sin información auxiliar de nubes disponible en la escena.

\section{Región de Interés (ROI) y Caso de Estudio}
El área de estudio se enfoca sobre el sector central y norte del Parque Nacional Blue Mountains, como se muestra en la Figura \ref{fig:parque_contexto}, colindante con el sector de Gospers Mountain y el cañón del río Grose en Nueva Gales del Sur. La región se caracteriza por mesetas elevadas de arenisca de Hawkesbury, profundos valles labrados y bosques densos de eucalipto seco y húmedo (\textit{Eucalyptus} spp.) con una carga de combustible excepcionalmente inflamable bajo temperaturas superiores a 40~$^\circ$C \cite{hamilton2020ecological}.

\begin{figure}[H]
    \centering
    \includegraphics[width=0.68\textwidth]{12. Area del Parque Blue Mountains.jpg}
    \caption{Delimitación del Parque Nacional Blue Mountains}
    \label{fig:parque_contexto}
\end{figure}

Dado que el parque cubre un área extensa y de geometría irregular, se decidió acotar el análisis a una Región de Interés (ROI) rectangular diagonal. Esta forma permite abarcar de manera continua el corredor por donde avanzó el frente de fuego del complejo Gospers Mountain, concentrando el gradiente de severidad sobre el dosel forestal y evitando incluir zonas no afectadas que aumentarían el costo computacional sin aportar información relevante.

\section{Fase 1: Definición del Problema y Adquisición de Datos}

\subsection{Criterio de Selección Temporal}
Para este análisis se trabajó con una única escena Sentinel-2, correspondiente a la condición posterior al incendio (19 de febrero de 2020). Esta fecha corresponde a la ventana inmediata posterior a la extinción del megaincendio: a inicios y mediados de febrero de 2020, precipitaciones torrenciales de gran escala apagaron definitivamente los frentes residuales y lavaron las partículas de humo de la columna atmosférica. Esto permite la visualización directa y descubierta del carbón, las cenizas y el dosel colapsado.

Al tratarse de una clasificación supervisada de una sola fecha, no se requiere una imagen previa al incendio: las coberturas de interés, incluida la vegetación boscosa quemada, se identifican directamente sobre esta escena.

\subsection{Plataforma y Ficha Técnica de las Escenas Satelitales}
Las imágenes fueron adquiridas a través del portal de acceso abierto \textit{Copernicus Data Space Ecosystem} (CDSE) de la Agencia Espacial Europea (ESA). A diferencia de misiones de resolución media como Landsat (30 m), la misión \textbf{Sentinel-2 (MSI)} ofrece bandas con tamaño de celda de 10 m y 20 m, permitiendo un nivel de detalle cartográfico superior.

\begin{table}[H]
    \centering
    \caption{Parámetros técnicos de las adquisiciones satelitales Sentinel-2 MSI.}
    \label{tab:metadatos_s2}
    \begin{tabular}{@{}ll@{}}
        \toprule
        \textbf{Parámetro} & \textbf{Especificación Técnica} \\
        \midrule
        Misión y Sensor & Sentinel-2 MSI (\textit{MultiSpectral Instrument}) \\
        Plataforma de Descarga & Copernicus Data (\url{https://dataspace.copernicus.eu/}) \\
        Nivel de Procesamiento & Level-2A (BOA - \textit{Bottom-Of-Atmosphere Surface Reflectance}) \\
        Teselas de la Grilla MGRS & \textbf{T56HKH} (sector central/sur) y \textbf{T56HKJ} (sector septentrional) \\
        Fecha de Adquisición (post-incendio) & 19 de febrero de 2020 \\
        Proyección de Trabajo & GDA2020 / MGA zone 56 (EPSG:7856) \\
        Bandas Utilizadas & 10 bandas: B01, B02, B03, B04, B05, B06, B07, B8A, B11 y B12 \\
        Resolución Espacial & Unificada a 20 m (remuestreo) \\
        Cobertura Nubosa en ROI & Cobertura marginal evaluada con capa de clasificación de escena \\
        \bottomrule
    \end{tabular}
\end{table}

\subsection{Necesidad Metodológica de Cobertura Multi-Tesela}
Al contrastar el área geográfica del parque Blue Mountains con la partición en teselas de la grilla, se observó que esta no puede ser cubierta por una única escena, como se ve en la Figura \ref{fig:seleccion_poligono_general}. Por lo tanto, se deben tomar dos divisiones, presentadas en la Figura \ref{fig:teselas_compartidas}, de esta forma la reserva queda dividida en la frontera entre la tesela inferior \textbf{T56HKH} y la tesela superior \textbf{T56HKJ}.

\begin{figure}[H]
    \centering
    \includegraphics[width=0.9\textwidth]{1. Seleccion Poligono Blue Mountains.jpg}
    \caption{Identificación del área de Blue Mountains sobre la plataforma Copernicus.}
    \label{fig:seleccion_poligono_general}
\end{figure}

\begin{figure}[H]
    \centering
    \includegraphics[width=0.9\textwidth]{2. Zona con teselas compartidas.jpg}
    \caption{Intercepción espacial del área de interés sobre dos teselas adyacentes.}
    \label{fig:teselas_compartidas}
\end{figure}

Por consiguiente, en el panel lateral de búsqueda de Copernicus se seleccionaron explícitamente ambas escenas (Figuras \ref{fig:panel_teselas} y \ref{fig:especificacion_teselas}) para la fecha evaluada, descargando los paquetes completos en formato SAFE Nivel-2A.

\begin{figure}[H]
    \centering
    \begin{minipage}{0.52\textwidth}
        \centering
        \includegraphics[width=\textwidth]{3. Selección de teselas en panel lateral.jpg}
        \caption{Selección de las escenas correspondientes a la zona de estudio.}
        \label{fig:panel_teselas}
    \end{minipage}\hfill
    \begin{minipage}{0.42\textwidth}
        \centering
        \includegraphics[width=\textwidth]{4. Especificacion de teselas.jpg}
        \caption{Inspección de metadatos de las teselas T56HKH y T56HKJ.}
        \label{fig:especificacion_teselas}
    \end{minipage}
\end{figure}

\section{Fase 2: Preprocesamiento Tabular y Entrenamiento de Modelos de Machine Learning}

\subsection{Construcción de la Base de Datos y Prevención de Fuga de Información (\textit{Data Leakage})}
A partir de los polígonos de entrenamiento y validación terrestre construidos sobre la escena Sentinel-2 de Blue Mountains, se extrajeron los valores de reflectancia de 10 bandas espectrales: aerosol costero (B01), azul (B02), verde (B03), rojo (B04), tres bandas de borde rojo (B05, B06, B07), infrarrojo cercano ancho (B08), e infrarrojos de onda corta (B11 y B12). Cada fila del conjunto tabulado representa la firma espectral puntual de un píxel asociada a una de las \textbf{cuatro macroclases} finales del proyecto: (1) Agua, (2) Urbanización, (3) Vegetación boscosa sana y (4) Vegetación boscosa quemada. Esta consolidación proviene de los polígonos originales de siete coberturas, sobre los cuales se aplicaron tres fusiones por recomendación del docente en la primera entrega: Pasto (ID~3) fue absorbido en Vegetación boscosa sana (ID~4) por similitud espectral; Nubes (ID~6) fue absorbido en Urbanización (ID~2) por alta reflectancia compartida en el visible; y Sombras (ID~7) fue eliminado por escasísima representatividad muestral.

El conjunto total consolidado comprende \textbf{70.519 píxeles muestreados}. Para garantizar la reproducibilidad y validez estadística de los experimentos, la base se particionó de manera previa en un \textbf{70\% para entrenamiento} ($n = 49.363$) y un \textbf{30\% independiente para prueba} ($n = 21.156$), aplicando un muestreo estratificado para conservar la proporción natural de clases. 

Para modelos sensibles a la escala (como Redes Neuronales o SVM), los datos fueron estandarizados ajustando la media y varianza exclusivamente sobre el conjunto de entrenamiento. Para los modelos basados en árboles de decisión se utilizaron los valores espectrales crudos (\textit{raw}), dado que los algoritmos de partición jerárquica son invariantes ante transformaciones monótonas de las variables predictoras \cite{breiman1984classification}, preservando la interpretabilidad física directa de los umbrales de reflectancia en las reglas de decisión.

Para asegurar un entorno reproducible y libre de sesgos entre los distintos modelos evaluados, el flujo tabular se estructuró mediante tres módulos automatizados en Python:
\begin{enumerate}
    \item \textbf{Extracción espacial y tabulación (\texttt{generar\_tabulaciones.py}):} Mediante operaciones vectorizadas en \texttt{NumPy} y algoritmos geométricos de inclusión puntual (\textit{ray-casting}), se intersectaron los polígonos de entrenamiento provenientes del archivo \texttt{AfterProy2.scpx} con el mosaico multibanda de 10 bandas a 20 m (\texttt{EPSG:7856}). Se filtraron las celdas exteriores (\texttt{NoData = 0}) y se eliminaron duplicados espaciales por par fila-columna, exportando el conjunto de entrenamiento (\texttt{AfterProy2\_Training\_Tabulado.tsv}) y el inventario completo de píxeles válidos de la escena para la fase de inferencia posterior.
    \item \textbf{Partición y estandarización (\texttt{base\_modelos.py}):} Sobre los 70.519 píxeles consolidados (Tabla \ref{tab:distribucion_clases}), se ejecutó una partición estratificada fija (\texttt{random\_state = 42}) en 70\% entrenamiento ($n = 49.363$) y 30\% prueba ($n = 21.156$). Posteriormente, se ajustó un transformador de puntuación estándar (\texttt{StandardScaler}, $\mu = 0, \sigma = 1$) exclusivamente sobre las bandas de entrenamiento, serializando el objeto (\texttt{scaler.pkl}) para garantizar coherencia radiométrica durante la evaluación y la futura predicción espacial.
    \item \textbf{Marco unificado de evaluación (\texttt{evaluar\_modelos.py}):} Se estandarizó el registro de la búsqueda de hiperparámetros (\texttt{historial\_hpo\_*.tsv}), el cálculo de métricas globales y categóricas (\texttt{metricas\_*.tsv}, \texttt{reporte\_clases\_*.tsv}) y la generación de matrices de confusión $4 \times 4$ a 300 dpi, asegurando comparabilidad directa entre los cinco clasificadores.
\end{enumerate}

\begin{table}[H]
    \centering
    \caption{Distribución muestral y representatividad espacial de las macroclases en el conjunto tabulado (resolución de celda: $20\text{ m} \times 20\text{ m}$).}
    \label{tab:distribucion_clases}
    \begin{tabular}{@{}clrrrr@{}}
        \toprule
        \textbf{ID} & \textbf{Macroclase} & \textbf{Píxeles Totales} & \textbf{Train (70\%)} & \textbf{Test (30\%)} & \textbf{Proporción (\%)} \\
        \midrule
        1 & Agua & 18.477 & 12.934 & 5.543 & 26,20\% \\
        2 & Urbanización$^*$ & 2.538 & 1.776 & 762 & 3,60\% \\
        4 & Veg. Boscosa Sana$^{**}$ & 5.010 & 3.507 & 1.503 & 7,10\% \\
        5 & Veg. Boscosa Quemada & 44.494 & 31.146 & 13.348 & 63,10\% \\
        \midrule
        \multicolumn{2}{@{}l}{\textbf{Total Consolidado}} & \textbf{70.519} & \textbf{49.363} & \textbf{21.156} & \textbf{100,00\%} \\
        \bottomrule
    \end{tabular}
    \vspace{2pt}
    \begin{flushleft}
    \footnotesize{$^*$ Incluye píxeles originalmente etiquetados como Nubes (ID~6), absorbidos por alta reflectancia compartida en el visible.}\\
    \footnotesize{$^{**}$ Incluye píxeles originalmente etiquetados como Pasto (ID~3), absorbidos por similitud espectral en imagen post-incendio.}\\
    \footnotesize{Nota: La clase Sombras (ID~7) fue eliminada por escasísima representatividad muestral (283 píxeles, 0,40\% del dataset original).}
    \end{flushleft}
\end{table}
\subsection{Regresión Logística (Logistic Regression - LR)}
La Regresión Logística multiclase se utilizó como modelo base (\textit{baseline}) de la comparación. Es un clasificador lineal: combina las 10 bandas espectrales estandarizadas mediante una suma ponderada y transforma el resultado en una probabilidad para cada una de las cuatro macroclases (función \textit{softmax}), asignando cada píxel a la clase de mayor probabilidad. Al ser lineal, genera fronteras de decisión planas en el espacio espectral, por lo que sirve como punto de referencia para medir cuánto aportan los modelos no lineales.

\subsubsection{Optimización de Hiperparámetros (HPO)}
El modelo se implementó con \texttt{LogisticRegression} de \texttt{scikit-learn}, entrenado sobre los datos estandarizados (\texttt{X\_train.csv}). Se ejecutó una búsqueda en cuadrícula (\texttt{GridSearchCV}) con validación cruzada estratificada de 3 pliegues, maximizando el $F_1\text{-score}$ macro sobre \textbf{32 combinaciones} (96 ajustes totales):

\begin{itemize}
    \item \textbf{Inverso de la regularización $L_2$ (\texttt{C}):} $[0{,}01;\ 0{,}1;\ 1;\ 10;\ 100;\ 1000;\ 10000;\ 100000]$. Valores pequeños implican una regularización más fuerte.
    \item \textbf{Algoritmo de optimización (\texttt{solver}):} \texttt{'lbfgs'} y \texttt{'newton-cg'}.
    \item \textbf{Ponderación de clases (\texttt{class\_weight}):} \texttt{None} vs. \texttt{'balanced'}.
\end{itemize}

No se incluyeron las penalizaciones $L_1$ ni \textit{ElasticNet} (solver \texttt{saga}), debido al alto costo computacional con aproximadamente 50.000 píxeles de entrenamiento.

La configuración óptima alcanzó un puntaje medio de validación cruzada de $F_1\text{-macro} = 0{,}9805 \pm 0{,}0012$, con \texttt{C = 100}, \texttt{solver = 'newton-cg'} y \texttt{class\_weight = None}. Se observó que una regularización débil (\texttt{C} alto) mejora el desempeño, lo que indica que el modelo necesita libertad para ajustar los coeficientes de las bandas.

\subsubsection{Evaluación del Rendimiento en el Conjunto de Test}
Al evaluar el modelo óptimo sobre el 30\% de datos independientes ($n =$ 21.156 píxeles), se obtuvo:
\begin{itemize}
    \item \textbf{Overall Accuracy (OA):} $0{,}9947$ ($99{,}47\%$).
    \item \textbf{Coeficiente Kappa de Cohen ($\kappa$):} $0{,}9899$.
    \item \textbf{Precisión Macro:} $0{,}9823$ | \textbf{Recall Macro:} $0{,}9844$ | \textbf{$F_1$-Score Macro:} $0{,}9833$.
    \item \textbf{$F_1$-Score Ponderado:} $0{,}9947$.
\end{itemize}

En la Tabla \ref{tab:reporte_lr} se detalla el comportamiento por macroclase en el conjunto de prueba.

\begin{table}[H]
    \centering
    \caption{Métricas detalladas por macroclase para el modelo Logistic Regression (LR) en el conjunto de test.}
    \label{tab:reporte_lr}
    \begin{tabular}{@{}lcccc@{}}
        \toprule
        \textbf{Macroclase} & \textbf{Precisión} & \textbf{Recall} & \textbf{$F_1$-Score} & \textbf{Soporte (Píxeles)} \\
        \midrule
        Agua & 1,0000 & 0,9991 & 0,9995 & 5.543 \\
        Urbanización & 0,9762 & 0,9698 & 0,9730 & 762 \\
        Veg. Boscosa Sana & 0,9549 & 0,9721 & 0,9634 & 1.503 \\
        Veg. Boscosa Quemada & 0,9980 & 0,9968 & 0,9974 & 13.348 \\
        \midrule
        \textbf{Promedio Macro} & \textbf{0,9823} & \textbf{0,9844} & \textbf{0,9833} & \textbf{21.156} \\
        \textbf{Promedio Ponderado} & \textbf{0,9947} & \textbf{0,9947} & \textbf{0,9947} & \textbf{21.156} \\
        \bottomrule
    \end{tabular}
\end{table}

Como muestra la matriz de confusión (Figura \ref{fig:cm_lr}), el modelo clasifica casi perfectamente las clases dominantes: el agua ($F_1 = 0{,}9995$) y la vegetación boscosa quemada ($F_1 = 0{,}9974$; 13.305 de 13.348 píxeles). Sus errores más importantes se concentran en la confusión entre \textbf{Vegetación boscosa quemada y Sana} (42 píxeles quemados clasificados como sanos, el mayor error relativo del modelo), y entre \textbf{Urbanización y Vegetación boscosa sana} (22 px urbanos asignados a bosque sano). Al ser lineal, la regresión logística no logra separar bien las clases con firmas espectrales gradualmente superpuestas, lo que justifica el uso de modelos no lineales como SVM o ANN.

\begin{figure}[H]
    \centering
    \includegraphics[width=0.72\textwidth]{matriz_confusion_lr.jpg}
    \caption{Matriz de confusión del modelo Logistic Regression (LR) sobre el conjunto de test (21.156 píxeles).}
    \label{fig:cm_lr}
\end{figure}
\subsection{Árboles de Decisión (Decision Trees - DT)}
Los árboles de decisión representan modelos lógicos de caja blanca que dividen de forma recursiva el espacio de características multidimensional en hiper-rectángulos ortogonales mediante umbrales espectrales óptimos \cite{breiman1984classification}.

\subsubsection{Optimización de Hiperparámetros (HPO)}
Para evitar el sobreajuste (\textit{overfitting}) intrínseco de los árboles profundos sobre datos satelitales, se ejecutó una búsqueda exhaustiva en cuadrícula (\texttt{GridSearchCV}) acoplada a una validación cruzada estratificada de 3 pliegues (\textit{3-fold Stratified CV}). La métrica objetivo de optimización fue el $F_1\text{-score}$ macro, asegurando igual ponderación para macroclases con cobertura espacial reducida. La malla exploró \textbf{720 combinaciones} en total:

\begin{itemize}
    \item \textbf{Criterio de impureza:} Gini vs. Entropía de Shannon vs. Log-loss.
    \item \textbf{Profundidad máxima (\texttt{max\_depth}):} $[5, 10, 15, 20, 25, \text{None}]$.
    \item \textbf{Muestras mínimas para división (\texttt{min\_samples\_split}):} $[2, 5, 10, 20]$.
    \item \textbf{Muestras mínimas por hoja (\texttt{min\_samples\_leaf}):} $[1, 2, 4, 8, 16]$.
    \item \textbf{Ponderación de clases (\texttt{class\_weight}):} $\text{None}$ vs. $\text{Balanced}$.
\end{itemize}

La configuración óptima encontrada correspondió a: \texttt{criterion = 'log\_loss'}, \texttt{max\_depth = 25}, \texttt{min\_samples\_split = 2}, \texttt{min\_samples\_leaf = 1} y \texttt{class\_weight = None}, con un puntaje de validación cruzada de $F_1\text{-macro} = 0{,}9859 \pm 0{,}0014$. El uso de log-loss como criterio de impureza maximiza la ganancia de información en la frontera de transición espectral entre áreas calcinadas y bosque intacto.

\subsubsection{Evaluación del Rendimiento en el Conjunto de Test}
El modelo optimizado se evaluó sobre el 30\% de datos independientes, alcanzando las siguientes métricas globales:
\begin{itemize}
    \item \textbf{Overall Accuracy (OA):} $0{,}9961$ ($99{,}61\%$).
    \item \textbf{Coeficiente Kappa de Cohen ($\kappa$):} $0{,}9926$, evidenciando una concordancia casi perfecta más allá de asignaciones aleatorias.
    \item \textbf{Precisión Macro:} $0{,}9875$ | \textbf{Recall Macro:} $0{,}9879$ | \textbf{$F_1$-Score Macro:} $0{,}9877$.
    \item \textbf{$F_1$-Score Ponderado:} $0{,}9961$.
\end{itemize}

En la Tabla \ref{tab:reporte_dt} se desglosa el desempeño por categoría de cobertura terrestre.

\begin{table}[H]
    \centering
    \caption{Métricas detalladas por macroclase para el modelo Decision Trees (DT) en el conjunto de test.}
    \label{tab:reporte_dt}
    \begin{tabular}{@{}lcccc@{}}
        \toprule
        \textbf{Macroclase} & \textbf{Precisión} & \textbf{Recall} & \textbf{$F_1$-Score} & \textbf{Soporte (Píxeles)} \\
        \midrule
        Agua & 0,9996 & 0,9989 & 0,9993 & 5.543 \\
        Urbanización & 0,9764 & 0,9790 & 0,9777 & 762 \\
        Veg. Boscosa Sana & 0,9760 & 0,9754 & 0,9757 & 1.503 \\
        Veg. Boscosa Quemada & 0,9981 & 0,9983 & 0,9982 & 13.348 \\
        \midrule
        \textbf{Promedio Macro} & \textbf{0,9875} & \textbf{0,9879} & \textbf{0,9877} & \textbf{21.156} \\
        \textbf{Promedio Ponderado} & \textbf{0,9961} & \textbf{0,9961} & \textbf{0,9961} & \textbf{21.156} \\
        \bottomrule
    \end{tabular}
\end{table}

La matriz de confusión resultante (Figura \ref{fig:cm_dt}) corrobora una discriminación sobresaliente. De los 13.348 píxeles de vegetación quemada, el árbol clasificó correctamente 13.325 ($99{,}83\%$). Las leves confusiones residuales se concentraron entre \textbf{Vegetación boscosa sana y Urbanización} (18 px en cada sentido) y entre \textbf{Vegetación quemada y Sana} (22 px quemados clasificados como sanos), errores atribuibles al solapamiento espectral en las longitudes de onda del infrarrojo.

\begin{figure}[H]
    \centering
    \includegraphics[width=0.72\textwidth]{matriz_confusion_dt.jpg}
    \caption{Matriz de confusión del modelo Decision Tree (DT) sobre el conjunto de test (21.156 píxeles).}
    \label{fig:cm_dt}
\end{figure}

\subsubsection{Análisis de Importancia de Bandas Espectrales (\textit{Feature Importance})}
Una de las mayores ventajas analíticas del árbol de decisión en teledetección es su capacidad inherente para revelar la influencia relativa de cada longitud de onda en la discriminación de coberturas. La impureza de corte acumulada permitió jerarquizar las bandas Sentinel-2, como se ilustra en la Figura \ref{fig:importancia_dt}.

\begin{figure}[H]
    \centering
    \includegraphics[width=0.82\textwidth]{importancia_bandas_dt.jpg}
    \caption{Importancia relativa de las bandas espectrales de Sentinel-2 estimada por el modelo Decision Tree.}
    \label{fig:importancia_dt}
\end{figure}

\begin{enumerate}
    \item \textbf{Banda 11 (SWIR-1 -- 1610 nm) con 59,87\% de importancia:} Se constituyó como el predictor dominante. La región del infrarrojo de onda corta es directamente sensible al contenido de humedad de la biomasa foliar y la acumulación de carbón pirogénico, variables diagnósticas fundamentales del megaincendio.
    \item \textbf{Banda 7 (Red Edge 3 -- 783 nm) con 23,75\% de importancia:} La región del borde rojo vegetativo captura la inflexión crítica entre la absorción fotoquímica de pigmentos y la dispersión por mesófilo celular, permitiendo segmentar el dosel arbóreo intacto frente a las copas chamuscadas.
    \item \textbf{Banda 1 (Aerosol costero -- 443 nm) con 6,36\% de importancia:} Contribuye a discriminar superficies reflectantes (urbanización, cuerpos de agua) de coberturas de vegetación densa.
\end{enumerate}

Las tres bandas líderes agrupan conjuntamente el \textbf{89,97\%} del peso decisorio global del modelo. En la Figura \ref{fig:arbol_estructura} se presenta la estructura de los primeros niveles jerárquicos del árbol de decisión, ilustrando los umbrales numéricos de reflectancia que rigen las primeras particiones.

\begin{figure}[H]
    \centering
    \includegraphics[width=0.98\textwidth]{estructura_arbol_dt.jpg}
    \caption{Diagrama jerárquico esquemático de los primeros tres niveles del árbol de decisión podado.}
    \label{fig:arbol_estructura}
\end{figure}

\subsection{Bono: Ensamble Random Forest (RF)}
Como contribución avanzada sugerida para superar las limitaciones de varianza del árbol individual, se implementó la arquitectura de ensamble \textbf{Random Forest (RF)} \cite{breiman2001random}. Random Forest construye una multiplicidad de árboles descorrelacionados mediante muestreo bootstrap con reemplazo de las observaciones y selección aleatoria de subespacios de bandas en cada división (\textit{feature bagging}).

\subsubsection{Optimización y Configuración de Hiperparámetros}
A través de \texttt{GridSearchCV} con 3 pliegues estratificados, se evaluaron 288 combinaciones de hiperparámetros (864 ajustes en total). La configuración óptima identificó un desempeño de validación cruzada de $F_1\text{-macro} = 0{,}9905 \pm 0{,}0011$ con:
\begin{itemize}
    \item \textbf{Número de estimadores (\texttt{n\_estimators}):} 100 árboles de decisión.
    \item \textbf{Profundidad máxima (\texttt{max\_depth}):} 15 niveles.
    \item \textbf{Criterio de selección de bandas (\texttt{max\_features}):} \texttt{'sqrt'} (selección de $\sqrt{10} \approx 3$ bandas aleatorias por nodo).
    \item \textbf{Muestras mínimas por hoja (\texttt{min\_samples\_leaf}):} 1.
    \item \textbf{Muestras mínimas por división (\texttt{min\_samples\_split}):} 5.
    \item \textbf{Balanceo de pesos (\texttt{class\_weight}):} \texttt{None}.
\end{itemize}

\subsubsection{Resultados y Rendimiento Comparativo}
El ensamble exhibió un incremento notable en la estabilidad y precisión de generalización respecto al árbol simple sobre el conjunto independiente ($n = 21.156$ píxeles):
\begin{itemize}
    \item \textbf{Overall Accuracy (OA):} $0{,}9974$ ($99{,}74\%$).
    \item \textbf{Coeficiente Kappa ($\kappa$):} $0{,}9951$.
    \item \textbf{Precisión Macro:} $0{,}9902$ | \textbf{Recall Macro:} $0{,}9929$.
    \item \textbf{$F_1$-Score Macro:} $0{,}9915$ (una ganancia neta respecto al árbol individual).
    \item \textbf{$F_1$-Score Ponderado:} $0{,}9974$.
\end{itemize}

En la Tabla \ref{tab:reporte_rf} se presenta el desglose por macroclase para Random Forest.

\begin{table}[H]
    \centering
    \caption{Métricas detalladas por macroclase para el ensamble Random Forest (RF) en el conjunto de test.}
    \label{tab:reporte_rf}
    \begin{tabular}{@{}lcccc@{}}
        \toprule
        \textbf{Macroclase} & \textbf{Precisión} & \textbf{Recall} & \textbf{$F_1$-Score} & \textbf{Soporte (Píxeles)} \\
        \midrule
        Agua & 1,0000 & 0,9989 & 0,9995 & 5.543 \\
        Urbanización & 0,9755 & 0,9948 & 0,9851 & 762 \\
        Veg. Boscosa Sana & 0,9866 & 0,9787 & 0,9826 & 1.503 \\
        Veg. Boscosa Quemada & 0,9988 & 0,9990 & 0,9989 & 13.348 \\
        \midrule
        \textbf{Promedio Macro} & \textbf{0,9902} & \textbf{0,9929} & \textbf{0,9915} & \textbf{21.156} \\
        \textbf{Promedio Ponderado} & \textbf{0,9974} & \textbf{0,9974} & \textbf{0,9974} & \textbf{21.156} \\
        \bottomrule
    \end{tabular}
\end{table}

En la Figura \ref{fig:cm_rf} se ilustra la matriz de confusión de Random Forest, donde se observa una discriminación casi perfecta: 13.335 aciertos de 13.348 en vegetación quemada, 1.471 aciertos en sana y 758 en urbanización.

\begin{figure}[H]
    \centering
    \includegraphics[width=0.72\textwidth]{matriz_confusion_rf.jpg}
    \caption{Matriz de confusión obtenida por el ensamble Random Forest (RF) en el conjunto de test (21.156 píxeles).}
    \label{fig:cm_rf}
\end{figure}

\subsubsection{Consistencia Espectroscópica en la Importancia de Bandas}
Al promediar la importancia de características a lo largo del ensamble, Random Forest atenúa el sesgo de partición única y produce un ranking espectral biofísicamente consistente (Figura \ref{fig:importancia_rf}):
\begin{itemize}
    \item \textbf{B12 (SWIR-2 -- 2190 nm):} $25{,}38\%$
    \item \textbf{B11 (SWIR-1 -- 1610 nm):} $19{,}58\%$
    \item \textbf{B08 (NIR -- 842 nm):} $16{,}97\%$
\end{itemize}

\begin{figure}[H]
    \centering
    \includegraphics[width=0.82\textwidth]{importancia_bandas_rf.jpg}
    \caption{Importancia de bandas espectrales estimada por el ensamble Random Forest (100 estimadores).}
    \label{fig:importancia_rf}
\end{figure}

\noindent \textbf{Interpretación Biofísica Teledetectada:} Las bandas B12, B11 y B08 concentran el \textbf{61{,}93\%} de la capacidad predictiva del ensamble. Esta convergencia valida cuantitativamente la base matemática del Índice de Área Quemada Normalizado (\textit{Normalized Burn Ratio}, NBR) \cite{key2006landscape}:
\begin{equation}
    \text{NBR} = \frac{\rho_{\text{NIR}} - \rho_{\text{SWIR}}}{\rho_{\text{NIR}} + \rho_{\text{SWIR}}} = \frac{\text{B08} - \text{B12}}{\text{B08} + \text{B12}}
\end{equation}
El algoritmo de ensamble identificó autónomamente que el contraste entre la alta reflectancia del dosel foliar sano en el NIR y la marcada reflectancia del sustrato calcinado y deshidratado en el SWIR conforma la señal radiométrica más diagnóstica de la catástrofe del Black Summer.

\subsection{Máquinas de Vectores de Soporte (Support Vector Machines - SVM)}
Las Máquinas de Vectores de Soporte (\textit{Support Vector Machines}, SVM) constituyen clasificadores de margen máximo que buscan construir hiperplanos óptimos de separación entre clases maximizando la distancia geométrica hacia las observaciones más cercanas o vectores de soporte. Mediante la aplicación del truco del núcleo (\textit{kernel trick}), el algoritmo proyecta implícitamente las 10 bandas espectrales estandarizadas hacia espacios de características de dimensión superior, permitiendo resolver fronteras radiométricas altamente no lineales con una sólida capacidad de generalización sobre coberturas complejas.

\subsubsection{Optimización de Hiperparámetros (HPO)}
El entrenamiento se implementó mediante el estimador \texttt{SVC} de \texttt{scikit-learn} sobre el conjunto estandarizado (\texttt{X\_train.csv}), utilizando una estrategia multiclase uno-contra-uno (\textit{one-vs-one}) y memoria caché ampliada (\texttt{cache\_size = 1000}). Se ejecutó una búsqueda estructurada en cuadrícula (\texttt{GridSearchCV}) con validación cruzada estratificada de 3 pliegues (\textit{3-fold Stratified CV}), maximizando el $F_1\text{-score}$ macro sobre \textbf{64 combinaciones} (192 ajustes totales). La configuración óptima identificada alcanzó un puntaje medio de validación cruzada de $F_1\text{-macro} = 0{,}9915 \pm 0{,}0017$, correspondiendo al núcleo gaussiano (\texttt{kernel = 'rbf'}) con penalización \texttt{C = 100}, parámetro \texttt{gamma = 0.1} y \texttt{class\_weight = None}.

\subsubsection{Evaluación del Rendimiento en el Conjunto de Test}
Al evaluar el modelo SVM óptimo sobre el 30\% de datos independientes ($n = 21.156$ píxeles), se obtuvieron las siguientes métricas globales:
\begin{itemize}
    \item \textbf{Overall Accuracy (OA):} $0{,}9978$ ($99{,}78\%$).
    \item \textbf{Coeficiente Kappa de Cohen ($\kappa$):} $0{,}9958$, reflejando una concordancia casi perfecta entre las predicciones y la verdad de terreno.
    \item \textbf{Precisión Macro:} $0{,}9901$ | \textbf{Recall Macro:} $0{,}9951$ | \textbf{$F_1$-Score Macro:} $0{,}9926$.
    \item \textbf{$F_1$-Score Ponderado:} $0{,}9978$.
\end{itemize}

En la Tabla \ref{tab:reporte_svm} se presenta el desglose de métricas por macroclase sobre el conjunto de prueba.

\begin{table}[H]
    \centering
    \caption{Métricas detalladas por macroclase para el modelo Support Vector Machine (SVM) en el conjunto de test.}
    \label{tab:reporte_svm}
    \begin{tabular}{@{}lcccc@{}}
        \toprule
        \textbf{Macroclase} & \textbf{Precisión} & \textbf{Recall} & \textbf{$F_1$-Score} & \textbf{Soporte (Píxeles)} \\
        \midrule
        Agua & 1,0000 & 0,9993 & 0,9996 & 5.543 \\
        Urbanización & 0,9769 & 0,9974 & 0,9870 & 762 \\
        Veg. Boscosa Sana & 0,9841 & 0,9854 & 0,9847 & 1.503 \\
        Veg. Boscosa Quemada & 0,9996 & 0,9986 & 0,9991 & 13.348 \\
        \midrule
        \textbf{Promedio Macro} & \textbf{0,9901} & \textbf{0,9951} & \textbf{0,9926} & \textbf{21.156} \\
        \textbf{Promedio Ponderado} & \textbf{0,9978} & \textbf{0,9978} & \textbf{0,9978} & \textbf{21.156} \\
        \bottomrule
    \end{tabular}
\end{table}

Como se observa en la matriz de confusión (Figura \ref{fig:cm_svm}), el clasificador SVM con núcleo RBF alcanzó una separación sobresaliente en las cuatro coberturas: 5.539 aciertos en agua ($99{,}93\%$), 760 aciertos en urbanización ($99{,}74\%$), 1.481 en vegetación sana ($98{,}54\%$) y 13.329 en vegetación quemada ($99{,}86\%$), reduciendo los errores de confusión cruzada a niveles prácticamente despreciables.

\begin{figure}[H]
    \centering
    \includegraphics[width=0.72\textwidth]{matriz_confusion_svm.jpg}
    \caption{Matriz de confusión del modelo Support Vector Machine (SVM) sobre el conjunto de test (21.156 píxeles).}
    \label{fig:cm_svm}
\end{figure}

\subsection{Redes Neuronales Artificiales (Artificial Neural Networks - ANN)}
Las Redes Neuronales Artificiales basadas en el Perceptrón Multicapa (\textit{Multi-Layer Perceptron}, MLP) constituyen aproximadores universales capaces de modelar fronteras de decisión altamente no lineales en espacios espectrales de alta dimensionalidad \cite{rumelhart1986learning}. A diferencia de los modelos basados en particiones ortogonales, la arquitectura neuronal proyecta las 10 bandas espectrales estandarizadas a través de capas ocultas densamente conectadas, capturando interacciones radiométricas complejas entre el visible (VIS), red edge y el SWIR.

\subsubsection{Optimización de Hiperparámetros (HPO)}
El entrenamiento se implementó mediante \texttt{MLPClassifier} utilizando el optimizador estocástico \texttt{Adam}, incorporando parada temprana (\texttt{early\_stopping = True}, con 10\% de validación interna y tolerancia de 10 épocas sin mejora) para prevenir el sobreajuste y optimizar el tiempo de convergencia. Se ejecutó una búsqueda sistemática en cuadrícula (\texttt{GridSearchCV}) con validación cruzada estratificada de 3 pliegues (\textit{3-fold Stratified CV}), maximizando el $F_1\text{-score}$ macro sobre \textbf{48 combinaciones} (144 ajustes totales). La configuración óptima identificada alcanzó un puntaje medio de validación cruzada de $F_1\text{-macro} = 0{,}9918 \pm 0{,}0013$, correspondiendo a una arquitectura profunda de dos capas ocultas con \textbf{64 y 32 neuronas} (\texttt{hidden\_layer\_sizes = (64, 32)}), función de activación tangente hiperbólica (\texttt{activation = 'tanh'}), regularización \texttt{alpha = 0.001} y tasa de aprendizaje inicial \texttt{learning\_rate\_init = 0.005}.

\subsubsection{Evaluación del Rendimiento en el Conjunto de Test}
Al evaluar la red neuronal óptima sobre el 30\% de datos independientes ($n = 21.156$ píxeles), el modelo obtuvo un desempeño global sobresaliente:
\begin{itemize}
    \item \textbf{Overall Accuracy (OA):} $0{,}9976$ ($99{,}76\%$).
    \item \textbf{Coeficiente Kappa de Cohen ($\kappa$):} $0{,}9954$.
    \item \textbf{Precisión Macro:} $0{,}9898$ | \textbf{Recall Macro:} $0{,}9946$ | \textbf{$F_1$-Score Macro:} $0{,}9922$.
    \item \textbf{$F_1$-Score Ponderado:} $0{,}9976$.
\end{itemize}

En la Tabla \ref{tab:reporte_ann} se detalla el comportamiento categórico por macroclase en el conjunto de prueba.

\begin{table}[H]
    \centering
    \caption{Métricas detalladas por macroclase para el modelo Artificial Neural Network (ANN) en el conjunto de test.}
    \label{tab:reporte_ann}
    \begin{tabular}{@{}lcccc@{}}
        \toprule
        \textbf{Macroclase} & \textbf{Precisión} & \textbf{Recall} & \textbf{$F_1$-Score} & \textbf{Soporte (Píxeles)} \\
        \midrule
        Agua & 1,0000 & 0,9993 & 0,9996 & 5.543 \\
        Urbanización & 0,9756 & 0,9987 & 0,9870 & 762 \\
        Veg. Boscosa Sana & 0,9840 & 0,9820 & 0,9830 & 1.503 \\
        Veg. Boscosa Quemada & 0,9994 & 0,9986 & 0,9990 & 13.348 \\
        \midrule
        \textbf{Promedio Macro} & \textbf{0,9898} & \textbf{0,9946} & \textbf{0,9922} & \textbf{21.156} \\
        \textbf{Promedio Ponderado} & \textbf{0,9976} & \textbf{0,9976} & \textbf{0,9976} & \textbf{21.156} \\
        \bottomrule
    \end{tabular}
\end{table}

Como evidencia la matriz de confusión (Figura \ref{fig:cm_ann}), la red neuronal demostró una notable precisión general: 5.539 aciertos en agua ($99{,}93\%$), 761 aciertos en urbanización ($99{,}87\%$), 1.476 en vegetación sana ($98{,}20\%$) y 13.329 en vegetación quemada ($99{,}86\%$), restringiendo las confusiones residuales a 19 píxeles entre vegetación sana y urbanización y 19 entre quemada y sana.

\begin{figure}[H]
    \centering
    \includegraphics[width=0.72\textwidth]{matriz_confusion_ann.jpg}
    \caption{Matriz de confusión del modelo Artificial Neural Network (ANN) sobre el conjunto de test (21.156 píxeles).}
    \label{fig:cm_ann}
\end{figure}

\subsection{K-Vecinos Más Cercanos (K-Nearest Neighbors - KNN)}
El algoritmo de K-Vecinos Más Cercanos es un método no paramétrico basado en instancias: no aprende una función explícita durante el entrenamiento, sino que almacena los píxeles de entrenamiento y clasifica cada nuevo píxel según la clase mayoritaria entre sus $k$ vecinos más próximos en el espacio espectral \cite{cover1967nearest}. Dado que la cercanía se mide mediante una distancia entre las 10 bandas, el modelo es sensible a la escala de las variables, por lo que se entrenó sobre los datos estandarizados (\texttt{X\_train.csv}), evitando que las bandas con mayores valores de reflectancia dominaran el cálculo de la distancia.

\subsubsection{Optimización de Hiperparámetros (HPO)}
El modelo se implementó con \texttt{KNeighborsClassifier} de \texttt{scikit-learn} \cite{pedregosa2011scikit}. Se ejecutó una búsqueda en cuadrícula (\texttt{GridSearchCV}) con validación cruzada estratificada de 3 pliegues, maximizando el $F_1\text{-score}$ macro sobre \textbf{54 combinaciones} (162 ajustes totales):

\begin{itemize}
    \item \textbf{Número de vecinos (\texttt{n\_neighbors}):} $[1, 3, 5, 7, 9, 11, 15, 21, 31]$. Valores pequeños producen fronteras más flexibles; valores grandes suavizan la decisión.
    \item \textbf{Ponderación del voto (\texttt{weights}):} \texttt{'uniform'} vs. \texttt{'distance'}.
    \item \textbf{Parámetro de la distancia de Minkowski (\texttt{p}):} $[1, 2, 3]$, correspondientes a las distancias Manhattan, Euclidiana y Minkowski de orden 3.
\end{itemize}

La configuración óptima alcanzó un puntaje medio de validación cruzada de $F_1\text{-macro} = 0{,}9890 \pm 0{,}0010$, con \texttt{n\_neighbors = 1}, \texttt{p = 1} (distancia Manhattan) y \texttt{weights = 'uniform'}. Con un solo vecino la ponderación por distancia produce resultados equivalentes.

\subsubsection{Evaluación del Rendimiento en el Conjunto de Test}
Al evaluar el modelo óptimo sobre el 30\% de datos independientes ($n = 21.156$ píxeles), se obtuvo:
\begin{itemize}
    \item \textbf{Overall Accuracy (OA):} $0{,}9966$ ($99{,}66\%$).
    \item \textbf{Coeficiente Kappa de Cohen ($\kappa$):} $0{,}9936$.
    \item \textbf{Precisión Macro:} $0{,}9868$ | \textbf{Recall Macro:} $0{,}9912$ | \textbf{$F_1$-Score Macro:} $0{,}9889$.
    \item \textbf{$F_1$-Score Ponderado:} $0{,}9967$.
\end{itemize}

En la Tabla \ref{tab:reporte_knn} se detalla el comportamiento por macroclase en el conjunto de prueba.

\begin{table}[H]
    \centering
    \caption{Métricas detalladas por macroclase para el modelo K-Nearest Neighbors (KNN) en el conjunto de test.}
    \label{tab:reporte_knn}
    \begin{tabular}{@{}lcccc@{}}
        \toprule
        \textbf{Macroclase} & \textbf{Precisión} & \textbf{Recall} & \textbf{$F_1$-Score} & \textbf{Soporte (Píxeles)} \\
        \midrule
        Agua & 1,0000 & 0,9993 & 0,9996 & 5.543 \\
        Urbanización & 0,9778 & 0,9843 & 0,9810 & 762 \\
        Veg. Boscosa Sana & 0,9698 & 0,9834 & 0,9765 & 1.503 \\
        Veg. Boscosa Quemada & 0,9994 & 0,9978 & 0,9986 & 13.348 \\
        \midrule
        \textbf{Promedio Macro} & \textbf{0,9868} & \textbf{0,9912} & \textbf{0,9889} & \textbf{21.156} \\
        \textbf{Promedio Ponderado} & \textbf{0,9967} & \textbf{0,9966} & \textbf{0,9967} & \textbf{21.156} \\
        \bottomrule
    \end{tabular}
\end{table}

Como muestra la matriz de confusión (Figura \ref{fig:cm_knn}), el modelo clasifica de manera muy consistente: 5.539 aciertos en agua ($99{,}93\%$), 750 en urbanización ($98{,}43\%$), 1.478 en vegetación sana ($98{,}34\%$) y 13.318 en vegetación quemada ($99{,}78\%$). Sus errores principales radican en 30 píxeles de vegetación quemada predichos como vegetación sana y 17 píxeles de vegetación sana asignados a urbanización.

\begin{figure}[H]
    \centering
    \includegraphics[width=0.72\textwidth]{matriz_confusion_knn.jpg}
    \caption{Matriz de confusión del modelo K-Nearest Neighbors (KNN) sobre el conjunto de test (21.156 píxeles).}
    \label{fig:cm_knn}
\end{figure}

\noindent \textbf{Consideraciones sobre el uso de KNN:} Que el valor óptimo sea $k = 1$ sugiere que los píxeles de una misma cobertura forman agrupaciones compactas en el espacio espectral. No obstante, al ser un método basado en instancias, KNN debe comparar cada nuevo píxel contra todo el conjunto de entrenamiento (49.363 píxeles), lo que lo hace significativamente más costoso computacionalmente en la fase de inferencia masiva sobre el mosaico raster.

\section{Análisis comparativo y elección del mejor modelo}

A partir de los resultados obtenidos sobre el conjunto de prueba independiente ($n = 21.156$ píxeles), se consolidó una evaluación comparativa de las cinco arquitecturas de Machine Learning implementadas (junto al ensamble Random Forest). La Tabla \ref{tab:comparativa_modelos} sintetiza el desempeño global de cada clasificador, ordenado jerárquicamente según su $F_1$-Score Macro.

\begin{table}[H]
    \centering
    \caption{Tabla comparativa de rendimiento global de los modelos de Machine Learning sobre el conjunto de test.}
    \label{tab:comparativa_modelos}
    \begin{tabular}{@{}lccccc@{}}
        \toprule
        \textbf{Modelo} & \textbf{Overall Accuracy} & \textbf{Kappa ($\kappa$)} & \textbf{$F_1$-Macro} & \textbf{$F_1$-Ponderado} & \textbf{Ranking} \\
        \midrule
        \textbf{SVM (RBF)} & \textbf{0,9978} & \textbf{0,9958} & \textbf{0,9926} & \textbf{0,9978} & \textbf{1} \\
        ANN (MLP) & 0,9976 & 0,9954 & 0,9922 & 0,9976 & 2 \\
        Random Forest (Bono) & 0,9974 & 0,9951 & 0,9915 & 0,9974 & 3 \\
        KNN & 0,9966 & 0,9936 & 0,9889 & 0,9967 & 4 \\
        Decision Tree (DT) & 0,9961 & 0,9926 & 0,9877 & 0,9961 & 5 \\
        Logistic Regression (LR) & 0,9947 & 0,9899 & 0,9833 & 0,9947 & 6 \\
        \bottomrule
    \end{tabular}
\end{table}

\subsection{Modelo Ganador - Support Vector Machine (SVM)}
El clasificador Support Vector Machine con núcleo de base radial (RBF) se consolida como la arquitectura óptima. Si bien todos los modelos superaron el 99\% de \textit{Overall Accuracy}, el análisis cuantitativo debe apoyarse en el $F_1$-Score Macro y el coeficiente Kappa, métricas rigurosas que ponderan equitativamente todas las coberturas independientemente de su abundancia muestral.

El SVM lideró ambas métricas ($F_1\text{-Macro} = 0{,}9926$ y $\kappa = 0{,}9958$), demostrando una superioridad estadística consistente. Su rendimiento sobresale en la delimitación de clases minoritarias con firmas espectrales intermedias: alcanzó un $F_1$-score de $0{,}9870$ en Urbanización (con un recall de $99{,}74\%$) y un $F_1$ de $0{,}9847$ en Vegetación boscosa sana. Por su parte, la red neuronal artificial (ANN) obtuvo un desempeño prácticamente idéntico ($F_1\text{-Macro} = 0{,}9922$), pero el SVM ofrece una garantía de convexidad global en su optimización cuadrática, minimizando el riesgo de estancamiento en mínimos locales que caracteriza a las arquitecturas neuronales. Por estas razones, SVM fue seleccionado como el modelo clasificador para la inferencia espacial del mosaico completo en la Fase 3.
\clearpage

\section{Fase 3: Clasificación del mosaico completo}
\label{sec:fase3}

\subsection{Predicción del ráster con el SVM}

En esta fase se aplicó el SVM entrenado al mosaico Sentinel-2 de Blue Mountains. Se desarrolló el script \texttt{predecir\_raster.py} para convertir las reflectancias de las diez bandas en un mapa categórico. El script toma como entrada \texttt{Mosaico\_Completo.tif} y carga el modelo, el escalador y el diccionario de clases guardados en el proyecto.

Como el mosaico tiene 5.247 columnas y 6.249 filas, el script lo lee con GDAL en bloques de hasta 128 filas para limitar el uso de memoria. En cada bloque conserva los píxeles con valores finitos y al menos una banda distinta de NoData. Para cada píxel válido, conserva el orden de las diez bandas, convierte los valores a \texttt{int32} como en la tabulación y aplica el escalador guardado en \texttt{scaler.pkl}. Con esos datos, el SVM predice una clase por píxel y el script escribe su código en la fila y columna correspondientes.

Para reproducir la predicción se coloca \texttt{Mosaico\_Completo.tif} junto a \texttt{predecir\_raster.py} y se ejecuta el script con el intérprete de Python de QGIS, que incluye GDAL.

\subsection{GeoTIFF generado y superficie por macroclase}

El script generó \texttt{Raster\_Clasificado.tif}, un GeoTIFF de una banda con los códigos de las cuatro macroclases. Se verificó que conserva las dimensiones, la transformación afín y el sistema de referencia del mosaico original, GDA2020 / MGA zona 56 (EPSG:7856). Cada celda mide $20\,\mathrm{m} \times 20\,\mathrm{m}$. El valor 0 quedó reservado para NoData.

Durante la escritura del GeoTIFF se contaron los píxeles asignados a cada código. Como cada celda cubre $400\,\mathrm{m}^2$ o $0{,}04\,\mathrm{ha}$, se multiplicó el conteo de cada clase por el área de una celda. El script guardó los conteos y las superficies en \texttt{areas\_por\_macroclase.csv}. La Tabla \ref{tab:fase3_superficie} muestra también la proporción de cada macroclase sobre los 16.071.574 píxeles válidos. Los conteos del CSV se verificaron contra los valores del GeoTIFF.

\begin{table}[H]
    \centering
    \caption{Superficie por macroclase en el área válida del mosaico clasificado.}
    \label{tab:fase3_superficie}
    \begin{tabular}{@{}clrrr@{}}
        \toprule
        \textbf{ID} & \textbf{Macroclase} & \textbf{Píxeles} & \textbf{Superficie (ha)} & \textbf{Área válida (\%)} \\
        \midrule
        1 & Agua & 509.115 & 20.364,60 & 3,17 \\
        2 & Urbanización & 1.937.647 & 77.505,88 & 12,06 \\
        4 & Vegetación boscosa sana & 5.615.829 & 224.633,16 & 34,94 \\
        5 & Vegetación boscosa quemada & 8.008.983 & 320.359,32 & 49,83 \\
        \midrule
        \multicolumn{2}{@{}l}{\textbf{Total}} & \textbf{16.071.574} & \textbf{642.862,96} & \textbf{100,00} \\
        \bottomrule
    \end{tabular}
\end{table}

La mayor superficie asignada por el modelo correspondió a Vegetación boscosa quemada, que cubre cerca de la mitad del área válida del mosaico.

\renewcommand{\refname}{Referencias}

\bibliographystyle{plain}
\bibliography{referencias}

\end{document}